% Generated draft: requires mathematical and source review.
\documentclass[11pt]{article}
\usepackage[margin=25mm]{geometry}
\usepackage{fontspec}
\setmainfont{Latin Modern Roman}
\usepackage{amsmath,amssymb,mathtools,unicode-math}
\setmathfont{Latin Modern Math}
\usepackage{enumitem,needspace,xurl,hyperref}
\setlength{\parindent}{0pt}
\setlength{\parskip}{5pt}
\setlength{\emergencystretch}{3em}
\newcommand{\N}{\mathbb{N}}
\newcommand{\Z}{\mathbb{Z}}
\newcommand{\Q}{\mathbb{Q}}
\newcommand{\R}{\mathbb{R}}
\newcommand{\C}{\mathbb{C}}
\newcommand{\F}{\mathbb{F}}
\newcommand{\A}{\mathbb{A}}
\newcommand{\PP}{\mathbb{P}}
\newcommand{\E}{\mathbb{E}}
\newcommand{\eps}{\varepsilon}
\newcommand{\dd}{\,\mathrm{d}}
\begin{document}
\section*{UnsolvedMath ID: 30000243}
\textbf{Source problem code:} catalogue entry
\subsection{Definitions and mathematical statement}
\textbf{Coefficient-image question.} The source describes stabilizing classifying spaces of elliptic bundles associated with genus-one conformal field theories to obtain generalized cohomology theories \(E\) equipped with natural maps \(\Phi:E\to K[[q]][q^{-1}]\). Are the images of their coefficient rings modular forms? The quantifier concerns theories and maps from this specified construction, not arbitrary cohomology theories or arbitrary maps into a Laurent-series target. The elliptic-bundle category, classifying spaces, stabilization, multiplicative structure, and maps are not constructed in the supplied record.\par In a conventional spectral interpretation, the coefficient ring is \(E_* = \pi_*(E)\), and the map induces \(\Phi_*:E_*\to K_*[[q]][q^{-1}]\), where \(K\) denotes the intended periodic \(K\)-theory and \(q\) is a formal variable. Here \([[q]][q^{-1}]\) allows Laurent series with only finitely many negative powers in each element. For complex periodic \(K\)-theory one writes \(K_*\cong\Z[\beta,\beta^{-1}]\), \(|\beta|=2\). After specifying the grading convention, an even-degree coefficient \(a\in E_{2w}\) has an image of the form \(\beta^w f_a(q)\). The usual modularity interpretation asks whether \(f_a(q)\) is the Fourier expansion at the chosen cusp of a modular form of the corresponding weight, for a specified modular group and coefficient ring. This is a conditional model of the claim; the source does not itself fix complex versus another \(K\)-theory, the weight-degree convention, or the normalization.\par A modular form of integer weight \(w\) for a specified subgroup \(\Gamma\) of \(\mathrm{SL}_2(\Z)\) is a holomorphic function on the upper half-plane satisfying \(f((a\tau+b)/(c\tau+d))=(c\tau+d)^w f(\tau)\) for \(\left(\begin{smallmatrix}a&b\\c&d\end{smallmatrix}\right)\in\Gamma\), with the prescribed behavior at cusps and any specified multiplier. Its Fourier expansion uses a cusp parameter, conventionally \(q=e^{2\pi i\tau}\) at a width-one cusp. Allowing poles at cusps gives a different, weakly holomorphic variant. Since the target allows negative powers of \(q\), the source must specify which cusp condition it intends; a Laurent target alone does not enforce holomorphic modular forms. The intended claim is image containment in the relevant Fourier-expansion ring, not surjectivity onto all modular forms.\par The supplied triage checked 2026-08-21 labels the item partially solved. It reports that elliptic spectra and topological modular forms realize major portions of the proposed picture, while the particular generalized-cohomology construction and coefficient-image claim are not a single fully settled theorem. No identification of this \(E\) with topological modular forms, or transfer of a theorem to every proposed construction, is inferred. This dated assessment requires independent verification.
\subsection{Short English statement}
For the cohomology theories arising from the specified elliptic-bundle construction, ask whether their coefficient images are Fourier expansions of modular forms. The construction, grading, modular group, and cusp conditions are missing; the question asks for containment, not recovery of every modular form.
\subsection{Sources}
\begin{itemize}\item Local UnsolvedMath registry, permanent ID 30000243, code catalogue entry; \emph{Modular Forms from Generalized Cohomology Theories}, Oberwolfach Reports Open Problems.\item \emph{Geometric Topology and Connections with Quantum Field Theory (2006)}, Oberwolfach Reports, DOI 10.4171/owr/2005/27, pp. 1505--1546. \url{https://doi.org/10.4171/owr/2005/27}. Proposed year and year tag: 2005; no proposer named.\item Supplied literature label \emph{Kriz--Sati generalized cohomology}: \url{https://arxiv.org/abs/hep-th/0410293}.\item Supplied literature label \emph{Elliptic cohomology overview}: \url{https://ncatlab.org/ncatlab/show/elliptic%20cohomology}. The supplied URL contains a TeX escape before the percent sign; the percent-encoded URL is rendered here.\item Supplied literature label \emph{OWR report}: \url{https://doi.org/10.4171/owr/2005/27}.\item Literature triage checked 2026-08-21, status \texttt{partially\_solved}; record updated 2026-09-16. The \texttt{corrected\_verified} field reports primary-report inspection, not independently checked here.\end{itemize}
\end{document}
